\documentclass[11pt]{article}
\usepackage[review]{acl}
\usepackage{times}
\usepackage{latexsym}
\usepackage{microtype}
\usepackage{inconsolata}

\title{ParabeLink: Retrieving Fables from Moral Lessons}

\author{Anonymous ACL submission}

\begin{document}
\maketitle

% ============================================================
%  ABSTRACT
% ============================================================
\begin{abstract}
Linking moral principles to the narratives that teach them poses a fundamental
retrieval challenge: morals and fables share virtually no vocabulary, demanding
conceptual understanding far beyond surface-level matching.
We introduce \textbf{ParabeLink}, a benchmark built on the MORABLES dataset
\citep{marcuzzo-etal-2025-morables} in which a system must identify the correct
fable for a given moral lesson from 709 candidates (mean lexical overlap IoU:
0.011).
The best off-the-shelf embedding model achieves MRR of 0.210, far below an oracle
ceiling of 0.893---computed by appending the ground-truth moral to each fable
document---a four-fold gap that surface matching cannot close.
Augmenting fable documents with LLM-generated moral-style summaries substantially
narrows this gap (MRR 0.360), demonstrating that the bottleneck lies in document
representation rather than encoder capacity.
Contrastive fine-tuning with moral-negative augmentation improves performance
further (MRR 0.438).
\textbf{ParabeLink} provides a challenging testbed for semantic retrieval where
success requires bridging a deep abstraction gap between query and document.
\end{abstract}

% ============================================================
%  SECTIONS (stubs — filled in session by session)
% ============================================================

\section{Introduction}
% TODO

\section{Related Work}

\subsection{Moral Reasoning and Figurative Language in NLP}

Moral and ethical reasoning poses a persistent challenge for NLP systems, as it requires abstracting principles from context rather than matching surface-level patterns.
Early efforts frame moral reasoning as classification: the ETHICS benchmark \citep{hendrycks2021aligning} evaluates language models on commonsense morality, deontological rules, and virtue across five ethical theory categories, while \citet{jiang2021delphi} introduce Delphi, a system trained to predict human moral judgments about everyday situations.
Both approaches demonstrate that neural models can approximate moral intuitions in constrained settings but struggle to generalize to novel ethical scenarios.
A related challenge arises in figurative and abstract language, where meaning departs fundamentally from surface form---idioms, metaphors, and parables require conceptual understanding that lexical overlap cannot supply \citep{TODO:idiolink}.
\citet{marcuzzo-etal-2025-morables} introduce MORABLES, the dataset on which \textbf{ParabeLink} is built, framing moral reasoning as multiple-choice selection over a small candidate set of Aesop's fables.
\textbf{ParabeLink} extends this challenge to open retrieval: given a moral, a system must identify the correct fable from 709 candidates with near-zero lexical overlap (IoU~=~0.011), demanding a level of abstraction that closed-set evaluation cannot expose.

\subsection{Semantic Retrieval and Embedding Models}

Neural dense retrieval---representing queries and documents as continuous vectors and retrieving by nearest-neighbor search---has become the dominant paradigm for open-domain information retrieval \citep{karpukhin2020dense}.
Pre-trained sentence encoders such as Sentence-BERT \citep{reimers2019sentence} and Contriever \citep{izacard2022unsupervised} demonstrated that contrastive objectives yield semantically meaningful representations without requiring labeled query-document pairs.
Subsequent work scaled this paradigm: E5 \citep{wang2022text} pre-trains on 270M weakly supervised text pairs; BGE \citep{xiao2023cpack} introduces a general-purpose embedding suite with hard-negative-aware contrastive fine-tuning; and Linq-Embed-Mistral \citep{kim2024linq} achieves state-of-the-art performance across the MTEB retrieval leaderboard.
Cross-domain generalization is evaluated by the BEIR benchmark \citep{thakur2021beir}, which shows that retrieval performance degrades substantially on out-of-domain tasks where surface cues are absent.
\textbf{ParabeLink} represents an extreme case along this axis: with near-zero lexical overlap between moral queries and fable documents, it isolates semantic abstraction as the primary bottleneck and directly stress-tests whether modern embedding models can reason across deep representational gaps.

\subsection{Contrastive Fine-tuning and Synthetic Data}

Contrastive learning is the dominant paradigm for training text embedding models: SimCSE \citep{gao2021simcse} demonstrates that dropout-based augmentation with in-batch negatives substantially improves representation quality, and hard negative mining---selecting negatives that are semantically plausible but factually distinct---further sharpens discriminative capacity.
In parallel, large language models have enabled synthetic data generation for retrieval, removing the need for costly annotation: InPars \citep{bonifacio2022inpars} uses few-shot prompting to generate synthetic queries for unlabeled documents, while subsequent work produces full query-positive-negative triplets for instruction-following retrieval \citep{wang2024improving}.
% TODO: fill in the 2-3 sentences describing exactly how we use STORAL augmentation and
% what distinguishes our contrastive fine-tuning setup from prior domain adaptation work
% (hint: the key claim is that our gap is *abstraction*, not just topical domain shift).
% Reference: \citep{TODO:storal}
Unlike prior domain adaptation work---where queries and documents share topical vocabulary---our setting requires bridging a deep abstraction gap; we show empirically that synthetic document augmentation narrows this gap substantially, and contrastive fine-tuning with moral-negative augmentation closes it further.

\section{The ParabeLink Benchmark}
% TODO

\section{Experiments}
% TODO

\section{Results and Discussion}
% TODO

\section{Conclusions}
% TODO

\section*{Limitations}
% TODO

\section*{Ethics Statement}
% TODO

% ============================================================
%  REFERENCES
% ============================================================
\bibliography{references}

\end{document}
